\section{Experimental Setup}
\label{sec:setup}

\draft{Setup section. Benchmarks follow the RLVR protocol of Sparrow
\citep{sadhukhan2026sparrow} so the two papers are measured on the same kind
of task. Everything here should be reproducible from the repo; cite the exact
scripts in the appendix.}

\paragraph{Models.}
Qwen3 thinking models \citep{yang2025qwen3} at 1.7B and 4B, and
\draft{8B if it fits on 48\,GB with the bounded cache}. \draft{Base-model
variants (Qwen3-4B-Base after a DeepScaleR warm-up, as in Sparrow's appendix)
are the cheaper option if 37k-token thinking rollouts are out of reach on
single GPUs.} The Qwen2.5 models from the memory sweeps
\citep{qwen2025qwen25} appear only in \cref{sec:memory} and \cref{app:memory}.

\paragraph{Training data.}
Mathematical reasoning: Polaris \citep{an2025polaris}, with DeepScaleR
\citep{luo2025deepscaler} as an optional warm-up stage to lengthen
chain-of-thought before the main run. Code: the 21K filtered TACO subset
\citep{li2023taco} with a sandboxed executor for the reward. All rewards are
binary and verifiable: final-answer match for math, unit-test pass for code.

\paragraph{Evaluation benchmarks.}
Math: AIME 2024, 2025 and 2026 (mean over 16 samples per problem), AMC 2023
and 2024 (mean over 4 samples), and MATH500 \citep{lightman2023letsverify}.
Code: LiveCodeBench \citep{jain2024livecodebench}, HumanEval+ and MBPP+
\citep{liu2023evalplus}. \draft{Pin the exact LiveCodeBench date window.}
Generation length cutoff \draft{TODO: 16k / 24k / 32k; Sparrow uses 37k for
thinking models, 12--24k for base models and code}.

\paragraph{Evaluation protocol.}
Every reported number is a \emph{paired} comparison on a fixed, seeded,
manifest-pinned prompt set (\cref{app:hygiene}). For each benchmark we report
accuracy (i) evaluated with dense attention and (ii) evaluated through the
same evicting cache used for training, so that the eviction gap and the
training gain are separated. We also report the training reward curve, which
is the primary stability signal in Sparrow and lets the two be read side by
side.

\paragraph{Baselines.}
\begin{itemize}
  \item \textbf{Dense-attention training} (GRPO, RLOO) wherever it fits in
    memory. This is the parity baseline.
  \item \textbf{Sparse-eval only}: base model evaluated through the cache, no
    training. This is the eviction gap.
  \item \textbf{SFT through the cache}: same data and cache, supervised
    objective on correct rollouts. Separates ``RL helps'' from ``any training
    helps''.
  \item \draft{Optional: sparse-rollout/dense-policy (Sparrow-style) at a
    length where dense gradient fits. Would make the mismatch comparison
    direct rather than argued. Defer the how-to-compare question.}
\end{itemize}

\paragraph{Hardware.}
Single-GPU throughout: NVIDIA L40S (48\,GB, g6e.2xlarge), run as a fleet of
independent workers. \draft{Report total GPU-hours per run.}

\paragraph{Metrics.}
Accuracy on the benchmarks above (avg@$k$ as stated), training reward, peak
allocated memory (\texttt{torch.cuda.max\_memory\_allocated}) and
CUDA-synchronised per-phase wall clock.
